\documentclass[10pt,a4paper]{amsart}
\usepackage{amsmath,amssymb,amsthm,graphicx}
% For commutative diagrams we use Xy-pic
\usepackage[all]{xy}

% We use 2cell for 2-commutative diagrams.
\xyoption{2cell}
% Theorem environments.
%
\theoremstyle{plain}
\newtheorem{theorem}[subsection]{Theorem}
\newtheorem{proposition}[subsection]{Proposition}
\newtheorem{lemma}[subsection]{Lemma}

\theoremstyle{definition}
\newtheorem{definition}[subsection]{Definition}
\newtheorem{example}[subsection]{Example}
\newtheorem{exercise}[subsection]{Exercise}
\newtheorem{situation}[subsection]{Situation}

\theoremstyle{remark}
\newtheorem{remark}[subsection]{Remark}
\newtheorem{remarks}[subsection]{Remarks}

\numberwithin{equation}{subsection}

% Macros
%
\def\lim{\mathop{\mathrm{lim}}\nolimits}
\def\colim{\mathop{\mathrm{colim}}\nolimits}
\def\Spec{\mathop{\mathrm{Spec}}}
\def\Hom{\mathop{\mathrm{Hom}}\nolimits}
\def\Ext{\mathop{\mathrm{Ext}}\nolimits}
\def\SheafHom{\mathop{\mathcal{H}\!\mathit{om}}\nolimits}
\def\SheafExt{\mathop{\mathcal{E}\!\mathit{xt}}\nolimits}
\def\Sch{\mathit{Sch}}
\def\Mor{\mathop{\mathrm{Mor}}\nolimits}
\def\Ob{\mathop{\mathrm{Ob}}\nolimits}
\def\Sh{\mathop{\mathit{Sh}}\nolimits}
\def\NL{\mathop{N\!L}\nolimits}
\def\CH{\mathop{\mathrm{CH}}\nolimits}
\def\proetale{{pro\text{-}\acute{e}tale}}
\def\etale{{\acute{e}tale}}
\def\QCoh{\mathit{QCoh}}
\def\Ker{\mathop{\mathrm{Ker}}}
\def\Im{\mathop{\mathrm{Im}}}
\def\Coker{\mathop{\mathrm{Coker}}}
\def\Coim{\mathop{\mathrm{Coim}}}

% Boxtimes
%
\DeclareMathSymbol{\boxtimes}{\mathbin}{AMSa}{"02}

%
% Macros for moduli stacks/spaces
%
\def\QCohstack{\mathcal{QC}\!\mathit{oh}}
\def\Cohstack{\mathcal{C}\!\mathit{oh}}
\def\Spacesstack{\mathcal{S}\!\mathit{paces}}
\def\Quotfunctor{\mathrm{Quot}}
\def\Hilbfunctor{\mathrm{Hilb}}
\def\Curvesstack{\mathcal{C}\!\mathit{urves}}
\def\Polarizedstack{\mathcal{P}\!\mathit{olarized}}
\def\Complexesstack{\mathcal{C}\!\mathit{omplexes}}
% \Pic is the operator that assigns to X its picard group, usage \Pic(X)
% \Picardstack_{X/B} denotes the Picard stack of X over B
% \Picardfunctor_{X/B} denotes the Picard functor of X over B
\def\Pic{\mathop{\mathrm{Pic}}\nolimits}
\def\Picardstack{\mathcal{P}\!\mathit{ic}}
\def\Picardfunctor{\mathrm{Pic}}
\def\Deformationcategory{\mathcal{D}\!\mathit{ef}}
\usepackage{cleveref}
\crefname{originallemma}{original lemma}{original lemmas}
\Crefname{originallemma}{Original lemma}{Original lemmas}
\makeatletter
\newlabel{flat-lemma-etale-flat-up-down}{{05B9}{source}{}{sourcepointer-1}{}}
\newlabel{flat-lemma-etale-flat-up-down@cref}{{[originallemma][0][]05B9}{[1][0][]source}{}{}{}}
\makeatother
\crefname{originalreference}{original reference}{original references}
\Crefname{originalreference}{Original reference}{Original references}
\makeatletter
\newlabel{lemma-sections-upstairs}{{0A3H}{source}{}{sourcepointer-2}{}}
\newlabel{lemma-sections-upstairs@cref}{{[originalreference][0][]0A3H}{[1][0][]source}{}{}{}}
\makeatother
\crefname{originallemma}{original lemma}{original lemmas}
\Crefname{originallemma}{Original lemma}{Original lemmas}
\makeatletter
\newlabel{sites-lemma-localize-morphism-strong}{{03EF}{source}{}{sourcepointer-3}{}}
\newlabel{sites-lemma-localize-morphism-strong@cref}{{[originallemma][0][]03EF}{[1][0][]source}{}{}{}}
\makeatother
\title{Base change for the pushforward of a sheaf of sets along a flat morphism with geometrically connected quasi-compact fibres over an etale base}
\author{The Stacks Project Authors}
\date{stacks-project:etale-cohomology.tex:local-snapshot-20261011; GNU FDL 1.2 or later}
\begin{document}
\maketitle
\noindent\textit{Editorial note.} All delivered mathematical text is the author's own, copied from the pinned Stacks Project chapter file etale-cohomology.tex; no statement or proof step was rewritten, reordered or removed.\par\medskip
\noindent\textit{Editorial note.} The counted proof is the author's complete argument: the proof is written out in full: the reduction to a good object with the required factorisation, the five-term chain of equalities with the reason for each, and the final passage from local agreement to the identity of sheaves.\par\medskip
\noindent\textit{Editorial note.} The label lemma-base-change-f-star-general names the counted statement, and the proof environment immediately following it in the source is delivered verbatim as its proof.\par\medskip
\noindent\textit{Editorial note.} The counted statement is carried by the author's own lemma environment in the chapter source.\par\medskip
\noindent\textit{Editorial note.} The statement and proof are copied from the pinned chapter file, including the author's footnote about locality of the pullback; the three cited inputs are kept as external pointers with their canonical Stacks tags.\par\medskip
\noindent\textit{Editorial note.} The preamble is the chapter's own native preamble excerpt already used for this source: the Xy-pic setup (native preamble lines 14-18) and the native theorem declarations with the mathematical macros (native preamble lines 154-217).\par\medskip
\noindent\textit{Allowed standard prerequisites (editorial).}\begin{itemize}
\item Sections upstairs for an etale morphism with geometrically connected fibres. Same chapter, the lemma computing the pullback of a pushforward on an object U whose structure morphism U -> V is quasi-compact with geometrically connected fibres. Supplies the first and third equalities in the displayed chain.
\item Flatness of an etale morphism that is flat up and down. More on Flatness, the lemma ensuring that U -> V is flat under the factorisation hypothesis. Supplies flatness of U -> V, one of the inputs to the sections-upstairs lemma.
\item Locality of the pullback along a morphism of sites. Sites, the lemma identifying the restriction of f\textasciicircum{}\{-1\}g\_*F to U\_etale with the pullback along U -> V of the restriction of g\_*F. Supplies the reduction of the first equality to the localised situation.
\end{itemize}\medskip
\ifdefined\hypertarget\hypertarget{sourcepointer-1}{}\fi
\noindent\textit{Source reference.} Original lemma 05B9: sources/stacks-chow/tags.txt, source lines 6877--6877. Prerequisite pointer invoked at proof line(s) 16631; it names a standard fact of a sibling chapter that the counted argument uses, and every step of the counted proof built on it is retained verbatim in the delivered body.\par\medskip
\ifdefined\hypertarget\hypertarget{sourcepointer-2}{}\fi
\noindent\textit{Source reference.} Original reference 0A3H: sources/stacks-etale-cohomology/pinned-original-source.tex, source lines 5030--5030. Prerequisite pointer invoked at proof line(s) 16642; it names a standard fact of the same chapter that the counted argument uses, and every step of the counted proof built on it is retained verbatim in the delivered body.\par\medskip
\ifdefined\hypertarget\hypertarget{sourcepointer-3}{}\fi
\noindent\textit{Source reference.} Original lemma 03EF: sources/stacks-chow/tags.txt, source lines 4454--4454. Prerequisite pointer invoked at proof line(s) 16646; it names a standard fact of a sibling chapter that the counted argument uses, and every step of the counted proof built on it is retained verbatim in the delivered body.\par\medskip
\begin{lemma}
\label{lemma-base-change-f-star-general}
Consider the cartesian diagram of schemes
$$
\xymatrix{
X \ar[d]_f & Y \ar[l]^h \ar[d]^e \\
S & T \ar[l]_g
}
$$
Assume that $f$ is flat and every object $U$ of $X_\etale$ has
a covering $\{U_i \to U\}$ such that $U_i \to S$
factors as $U_i \to V_i \to S$ with $V_i \to S$
\'etale and $U_i \to V_i$ quasi-compact with
geometrically connected fibres.
Then for any sheaf $\mathcal{F}$ of sets on $T_\etale$ we have
$f^{-1}g_*\mathcal{F} = h_*e^{-1}\mathcal{F}$.
\end{lemma}

\begin{proof}
Let $U \to X$ be an \'etale morphism such that $U \to S$ factors as
$U \to V \to S$ with $V \to S$ \'etale and $U \to V$ quasi-compact
with geometrically connected fibres. Observe that $U \to V$ is flat
(More on Flatness, Lemma \ref{flat-lemma-etale-flat-up-down}).
We claim that
\begin{align*}
f^{-1}g_*\mathcal{F}(U)
& = g_*\mathcal{F}(V) \\
& = \mathcal{F}(V \times_S T) \\
& = e^{-1}\mathcal{F}(U \times_X Y) \\
& = h_*e^{-1}\mathcal{F}(U)
\end{align*}
Namely, thinking of $U$ as an object of $X_\etale$ and
$V$ as an object of $S_\etale$ we see that the first equality
follows from Lemma \ref{lemma-sections-upstairs}\footnote{Strictly
speaking, we are also using that the restriction of $f^{-1}g_*\mathcal{F}$
to $U_\etale$ is the pullback via $U \to V$ of the restriction of
$g_*\mathcal{F}$ to $V_\etale$. See
Sites, Lemma \ref{sites-lemma-localize-morphism-strong}.}.
Thinking of $V \times_S T$ as an object of $T_\etale$
the second equality follows from the definition of $g_*$.
Observe that $U \times_X Y = U \times_S T$ (because $Y = X \times_S T$)
and hence $U \times_X Y \to V \times_S T$
has geometrically connected fibres as a base change of $U \to V$.
Thinking of $U \times_X Y$ as an object of $Y_\etale$, we see that
the third equality  follows from Lemma \ref{lemma-sections-upstairs}
as before. Finally, the fourth equality follows from the definition
of $h_*$.

\medskip\noindent
Since by assumption every object of $X_\etale$ has an \'etale
covering to which the argument of the previous paragraph applies
we see that the lemma is true.
\end{proof}

\end{document}
